% ============================================================
% GUÍA MAESTRA DE ESTUDIO — ÁRBOLES DE DECISIÓN Y ENSAMBLE
% Emanuel Cabrera Novoa · Análisis Predictivo de Datos · 2026
% Compilar con: pdflatex (2 pasadas) o lualatex
% ============================================================
\documentclass[11pt, a4paper]{book}

% ── Paquetes ─────────────────────────────────────────────────
\usepackage[utf8]{inputenc}
\usepackage[T1]{fontenc}
\usepackage[spanish]{babel}
\usepackage{lmodern}
\usepackage{microtype}
\usepackage[margin=2.2cm, top=2.5cm, bottom=2.5cm]{geometry}

% Colores y cajas
\usepackage{xcolor}
\usepackage{tcolorbox}
\tcbuselibrary{breakable, skins, theorems}

% Matemáticas
\usepackage{amsmath, amssymb, mathtools}

% Código fuente
\usepackage{listings}
\usepackage{inconsolata}

% Tablas
\usepackage{booktabs}
\usepackage{longtable}
\usepackage{array}
\usepackage{multirow}

% Gráficos y diagramas
\usepackage{tikz}
\usetikzlibrary{shapes.geometric, arrows.meta, positioning, fit, calc,
                decorations.pathreplacing, backgrounds}

% Otros
\usepackage{enumitem}
\usepackage{hyperref}
\usepackage{graphicx}
\usepackage{fancyhdr}
\usepackage{titlesec}
\usepackage{tocloft}
\usepackage{fontawesome5}

% ── Paleta de colores ─────────────────────────────────────────
\definecolor{primary}{HTML}{1A237E}      % azul oscuro
\definecolor{secondary}{HTML}{283593}
\definecolor{accent}{HTML}{E53935}       % rojo
\definecolor{treegreen}{HTML}{2E7D32}
\definecolor{baggingblue}{HTML}{1565C0}
\definecolor{boostingorange}{HTML}{E65100}
\definecolor{xgbpurple}{HTML}{6A1B9A}
\definecolor{lightbg}{HTML}{F5F5F5}
\definecolor{warningbg}{HTML}{FFF8E1}
\definecolor{successbg}{HTML}{E8F5E9}
\definecolor{infobg}{HTML}{E3F2FD}
\definecolor{critbg}{HTML}{FFEBEE}
\definecolor{codebg}{HTML}{263238}
\definecolor{codetext}{HTML}{ECEFF1}

% ── Estilo de página ─────────────────────────────────────────
\pagestyle{fancy}
\fancyhf{}
\fancyhead[LE,RO]{\small\thepage}
\fancyhead[LO]{\small\leftmark}
\fancyhead[RE]{\small Árboles de Decisión y Métodos de Ensamble}
\renewcommand{\headrulewidth}{0.4pt}

% ── Estilo de títulos ─────────────────────────────────────────
\titleformat{\chapter}[display]
  {\normalfont\huge\bfseries\color{primary}}
  {\chaptertitlename\ \thechapter}{10pt}{\Huge}
\titleformat{\section}
  {\normalfont\Large\bfseries\color{secondary}}
  {\thesection}{1em}{}
\titleformat{\subsection}
  {\normalfont\large\bfseries\color{secondary!80!black}}
  {\thesubsection}{1em}{}
\titleformat{\subsubsection}
  {\normalfont\normalsize\bfseries\color{secondary!60!black}}
  {\thesubsubsection}{1em}{}

% ── Cajas de colores ─────────────────────────────────────────
\tcbset{
  base/.style={
    breakable, enhanced,
    fonttitle=\bfseries\small,
    coltitle=white,
    attach boxed title to top left={yshift=-2mm, xshift=5mm},
    boxed title style={rounded corners, size=small},
    top=4mm
  }
}

\newtcolorbox{critico}[1][]{base,
  colback=critbg, colframe=accent, title={\faExclamationTriangle\ CONCEPTO CRÍTICO #1}}

\newtcolorbox{cuando}[1][]{base,
  colback=infobg, colframe=baggingblue, title={\faQuestion\ ¿CUÁNDO USAR? #1}}

\newtcolorbox{tip}[1][]{base,
  colback=successbg, colframe=treegreen, title={\faLightbulb\ TIP DE ESTUDIO #1}}

\newtcolorbox{alerta}[1][]{base,
  colback=warningbg, colframe=boostingorange, title={\faExclamationTriangle\ ADVERTENCIA #1}}

\newtcolorbox{formula}[1][]{base,
  colback=lightbg, colframe=primary, title={\faSuperscript\ FÓRMULA #1}}

\newtcolorbox{ejercicio}[1][]{base,
  colback=successbg!50, colframe=treegreen!60!black, title={\faPen\ EJERCICIO #1}}

\newtcolorbox{resumen}[1][]{base,
  colback=primary!8, colframe=primary!60, title={\faBookOpen\ RESUMEN #1}}

% ── Estilo de código ─────────────────────────────────────────
\lstdefinestyle{python}{
  language=Python,
  basicstyle=\ttfamily\footnotesize\color{codetext},
  backgroundcolor=\color{codebg},
  keywordstyle=\color{cyan}\bfseries,
  stringstyle=\color{yellow!70!orange},
  commentstyle=\color{green!60!white}\itshape,
  numberstyle=\tiny\color{gray},
  numbers=left,
  numbersep=8pt,
  breaklines=true,
  showstringspaces=false,
  frame=single,
  framesep=4pt,
  rulecolor=\color{codebg!50!black},
  tabsize=4,
  morekeywords={import, from, as, True, False, None, print, fit, predict,
                transform, fit_transform, score, shape},
}
\lstset{style=python}

% ── Hyperref ─────────────────────────────────────────────────
\hypersetup{
  colorlinks=true,
  linkcolor=primary,
  urlcolor=baggingblue,
  citecolor=treegreen,
  pdftitle={Guía Maestra - Árboles de Decisión y Ensamble},
  pdfauthor={Emanuel Cabrera Novoa},
}

% ────────────────────────────────────────────────────────────
\begin{document}

% ── PORTADA ──────────────────────────────────────────────────
\begin{titlepage}
  \begin{tikzpicture}[remember picture, overlay]
    \fill[primary] (current page.north west) rectangle
        ([yshift=-8cm]current page.north east);
    \fill[accent]  ([yshift=-8cm]current page.north west) rectangle
        ([yshift=-8.6cm]current page.north east);
  \end{tikzpicture}

  \vspace*{1.5cm}
  {\color{white}\fontsize{36}{44}\selectfont\bfseries
    Árboles de Decisión\\[4pt]
    y Métodos de Ensamble}\\[6pt]
  {\color{white!80!gray}\Large Guía Maestra de Estudio — Handbook Técnico Profesional}

  \vspace{3cm}
  \begin{tikzpicture}
    \node[fill=lightbg, rounded corners=8pt, inner sep=14pt, text width=13cm]{
      \begin{tabular}{ll}
        \textbf{Curso} & Análisis Predictivo de Datos · Machine Learning · Período 6\\[4pt]
        \textbf{Estudiante} & Emanuel Cabrera Novoa\\[4pt]
        \textbf{Docente} & Jhon Quiza\\[4pt]
        \textbf{Año} & 2026\\[4pt]
        \textbf{Versión} & 2.0 — Edición ampliada con guías de ejercicios
      \end{tabular}
    };
  \end{tikzpicture}

  \vfill
  \textit{Nota: compilar con \texttt{pdflatex} (dos pasadas) o \texttt{lualatex}.}
\end{titlepage}

% ── TABLA DE CONTENIDOS ──────────────────────────────────────
\tableofcontents
\newpage

% ════════════════════════════════════════════════════════════
\chapter{Flujo de Decisión — ¿Qué Modelo Usar?}
% ════════════════════════════════════════════════════════════

Esta sección es el mapa de navegación del curso. Antes de escribir una sola línea de código,
recorre este flujo para elegir el algoritmo adecuado.

\section{Diagrama de Flujo General}

\begin{center}
\begin{tikzpicture}[
  node distance=1.4cm and 2.2cm,
  decision/.style={diamond, draw=primary, fill=infobg, thick,
                   text width=3.2cm, align=center, inner sep=2pt, font=\footnotesize},
  process/.style={rectangle, rounded corners=6pt, draw=treegreen, fill=successbg,
                  thick, text width=3.8cm, align=center, font=\footnotesize\bfseries,
                  inner sep=6pt},
  terminal/.style={rectangle, rounded corners=12pt, draw=accent, fill=critbg,
                   thick, text width=4cm, align=center, font=\small\bfseries,
                   inner sep=7pt},
  arrow/.style={-Stealth, thick, color=secondary},
  label/.style={font=\scriptsize\bfseries, color=secondary}
]

% Nodo inicio
\node[terminal] (start) {TENGO UN PROBLEMA\\DE MACHINE LEARNING};

% Pregunta 1: ¿Interpretabilidad?
\node[decision, below=of start] (q1) {¿Necesito\\interpretabilidad\\total?};

% Árbol podado
\node[process, right=2.5cm of q1] (tree) {\faTree\ Árbol de Decisión\\Podado\\max\_depth=3--5};

% Pregunta 2: ¿Tamaño?
\node[decision, below=of q1] (q2) {¿Datos muy\\grandes?\\($>$500k filas)};

% LightGBM
\node[process, right=2.5cm of q2] (lgbm) {\faRocket\ LightGBM\\(histogramas+GOSS)};

% Pregunta 3: ¿Muchas categóricas?
\node[decision, below=of q2] (q3) {¿Muchas variables\\categóricas de alta\\cardinalidad?};

% CatBoost
\node[process, right=2.5cm of q3] (cat) {\faCat\ CatBoost\\(Ordered TS)};

% Pregunta 4: ¿Prototipo rápido?
\node[decision, below=of q3] (q4) {¿Prototipo rápido\\/ baseline?};

% Random Forest
\node[process, right=2.5cm of q4] (rf) {\faTree\ Random Forest\\(robusto, poco tuning)};

% Pregunta 5: ¿Máxima precisión?
\node[decision, below=of q4] (q5) {¿Datos ruidosos\\u outliers?};

% GB lr bajo
\node[process, right=2.5cm of q5] (gb) {\faChartLine\ Gradient Boosting\\lr bajo (0.01--0.05)};

% XGBoost (default para máx. precisión)
\node[process, below=of q5] (xgb) {\faBolt\ XGBoost\\(máxima precisión)};

% Flechas
\draw[arrow] (start) -- (q1);
\draw[arrow] (q1) -- node[label, above]{Sí} (tree);
\draw[arrow] (q1) -- node[label, right]{No} (q2);
\draw[arrow] (q2) -- node[label, above]{Sí} (lgbm);
\draw[arrow] (q2) -- node[label, right]{No} (q3);
\draw[arrow] (q3) -- node[label, above]{Sí} (cat);
\draw[arrow] (q3) -- node[label, right]{No} (q4);
\draw[arrow] (q4) -- node[label, above]{Sí} (rf);
\draw[arrow] (q4) -- node[label, right]{No} (q5);
\draw[arrow] (q5) -- node[label, above]{Sí} (gb);
\draw[arrow] (q5) -- node[label, right]{No} (xgb);
\end{tikzpicture}
\end{center}

\section{Flujo de Regularización de un Árbol}

\begin{center}
\begin{tikzpicture}[
  node distance=1.2cm,
  flowstep/.style={rectangle, rounded corners=5pt, draw=secondary, fill=infobg,
               thick, text width=10cm, align=left, inner sep=8pt, font=\footnotesize},
  check/.style={diamond, draw=accent, fill=critbg, thick,
                text width=3cm, align=center, font=\footnotesize},
  arrow/.style={-Stealth, thick, color=secondary}
]
\node[flowstep] (s1) {\textbf{PASO 1.} Entrenar árbol sin restricciones (\texttt{DecisionTreeClassifier()})\\
  Medir: Train acc y Test acc};
\node[check, below=0.7cm of s1] (c1) {Train $\approx$ 100\%\\Test $\ll$ Train?};
\node[flowstep, below=0.7cm of c1] (s2) {\textbf{PASO 2.} Aplicar pre-poda: probar \texttt{max\_depth} $\in \{3,4,5\}$, \texttt{min\_samples\_leaf} $\in \{3,5,7\}$\\
  Usar \texttt{GridSearchCV} con \texttt{scoring='f1\_weighted'}};
\node[flowstep, below=0.7cm of s2] (s3) {\textbf{PASO 3.} Aplicar post-poda: buscar \texttt{ccp\_alpha} en \texttt{np.logspace(-3, 3)}\\
  Comparar con pre-poda. Elegir el que generalice mejor.};
\node[flowstep, below=0.7cm of s3] (s4) {\textbf{PASO 4.} Verificar: Train acc $\approx$ Test acc (diferencia $<$ 5\%)\\
  Si sigue sobreajustando: reducir más la profundidad o aumentar \texttt{ccp\_alpha}};

\draw[arrow] (s1) -- (c1);
\draw[arrow] (c1) -- node[right, font=\scriptsize\bfseries, color=accent]{Sí = overfitting} (s2);
\draw[arrow] (s2) -- (s3);
\draw[arrow] (s3) -- (s4);
\end{tikzpicture}
\end{center}

\section{Flujo de Tuning para Boosting}

\begin{enumerate}[leftmargin=*, label=\textbf{Paso \arabic*.}]
  \item \textbf{Baseline rápido:} entrenar con parámetros por defecto. Obtener F1/accuracy base.
  \item \textbf{Número óptimo de árboles:} usar \textit{early stopping} con \texttt{learning\_rate=0.1}.
        Observar a qué iteración se estabiliza el error de validación.
  \item \textbf{Estructura del árbol:} variar \texttt{max\_depth} (3--8) y \texttt{min\_child\_weight}
        o \texttt{min\_samples\_leaf}. Más profundidad = menos sesgo, más varianza.
  \item \textbf{Regularización:} submustreo (\texttt{subsample=0.6--0.8}),
        columnas (\texttt{colsample\_bytree=0.6--0.8}),
        L1/L2 (\texttt{reg\_alpha}, \texttt{reg\_lambda}).
  \item \textbf{Ajuste fino del learning rate:} reducir a 0.05 o 0.01 y aumentar \texttt{n\_estimators}
        proporcionalmente. Re-validar con early stopping.
  \item \textbf{Desbalance de clases:} si recall de clase minoritaria es muy bajo, ajustar
        \texttt{scale\_pos\_weight} (XGBoost) o el umbral de decisión.
\end{enumerate}

\begin{cuando}[— Cuándo aplicar cada flujo]
  \begin{tabular}{lp{9cm}}
    \toprule
    \textbf{Situación} & \textbf{Flujo a seguir}\\
    \midrule
    Primera vez con el dataset & Flujo general (sección 1.1) para elegir familia\\
    Árbol único sobreajusta & Flujo de regularización (sección 1.2)\\
    Quiero la máxima precisión & Flujo de tuning boosting (sección 1.3)\\
    Presupuesto de tiempo limitado & Random Forest con \texttt{RandomizedSearchCV}, 30 iteraciones\\
    Competencia Kaggle & Stacking: RF + XGBoost + LightGBM + CatBoost\\
    \bottomrule
  \end{tabular}
\end{cuando}

% ════════════════════════════════════════════════════════════
\chapter{Árboles de Decisión — Teoría y Práctica}
% ════════════════════════════════════════════════════════════

\section{¿Qué es un Árbol de Decisión?}

Un árbol de decisión es una estructura jerárquica que divide el espacio de características mediante
reglas del tipo \textit{si–entonces}. Imita el razonamiento humano: una cadena de preguntas binarias
conduce a una decisión final.

\begin{resumen}[— Elementos del árbol]
  \begin{tabular}{lll}
    \toprule
    \textbf{Elemento} & \textbf{Descripción} & \textbf{Analogía}\\
    \midrule
    Nodo raíz (Root) & Primera pregunta, característica más importante &
      La pregunta más discriminativa\\
    Nodos interiores & Preguntas intermedias & Condiciones anidadas\\
    Ramas (Branches) & Posibles valores del atributo & Respuestas posibles\\
    Hojas (Leaf nodes) & Predicciones finales & Resultado: clase o valor\\
    \bottomrule
  \end{tabular}
\end{resumen}

\subsection{El Algoritmo CART}

CART (Classification and Regression Trees) construye árboles de forma \textit{greedy} y
\textit{top-down}: en cada nodo busca exhaustivamente el atributo y el punto de corte que
mejor divida los datos.

\begin{lstlisting}[caption={Pseudocódigo CART}]
def CART(nodo, datos, profundidad):
    if criterio_parada(datos, profundidad):
        return hoja(clase_mayoritaria(datos))
    
    mejor_ganancia = -inf
    for atributo a in caracteristicas:
        for valor de corte v in a:
            izq = datos[a <= v]
            der = datos[a > v]
            ganancia = impureza(datos) - impureza_ponderada(izq, der)
            if ganancia > mejor_ganancia:
                mejor = (a, v, ganancia)
    
    nodo.izq = CART(nodo_izq, izq, profundidad+1)
    nodo.der = CART(nodo_der, der, profundidad+1)
    return nodo
\end{lstlisting}

\begin{tip}[— Criterios de parada]
  El árbol deja de crecer cuando: (1) el nodo es puro (todos de la misma clase),
  (2) se alcanzó \texttt{max\_depth}, (3) el nodo tiene menos muestras que
  \texttt{min\_samples\_split}, o (4) la ganancia de información es cero.
\end{tip}

\section{Criterios de Impureza}

\subsection{Entropía}

Proviene de la teoría de la información de Shannon (1948). Mide la incertidumbre de un conjunto.

\begin{formula}[— Entropía de Shannon]
  \[
    H(S) = -\sum_{c \in \mathcal{C}} p(c) \cdot \log_2 p(c)
  \]
  \begin{itemize}
    \item $S$ = conjunto de datos, $\mathcal{C}$ = conjunto de clases, $p(c)$ = proporción de clase $c$
    \item $H(S) = 0$ $\Rightarrow$ nodo completamente puro (certeza máxima)
    \item $H(S) = 1$ $\Rightarrow$ máxima impureza para 2 clases (50/50)
    \item Rango: $\bigl[0,\; \log_2|\mathcal{C}|\bigr]$
  \end{itemize}
\end{formula}

\subsubsection{Ejemplo numérico — Dataset de Tenis}

Datos: 4 días se jugó tenis (Yes), 3 no se jugó (No).
\[
  p(\text{Yes}) = \tfrac{4}{7} \approx 0.571,\quad p(\text{No}) = \tfrac{3}{7} \approx 0.429
\]
\[
  H(S) = -(0.571 \cdot \log_2 0.571) - (0.429 \cdot \log_2 0.429)
       = 0.461 + 0.524 = \mathbf{0.985}
\]
\textit{Interpretación:} casi máxima impureza; los datos están muy mezclados.

\subsection{Ganancia de Información}

\begin{formula}[— Ganancia de Información]
  \[
    GI(S, a) = H(S) - \sum_{v \in \text{values}(a)} \frac{|S_v|}{|S|} \cdot H(S_v)
  \]
  El árbol selecciona el atributo con \textbf{mayor} ganancia (equivale al de \textbf{menor}
  entropía ponderada tras la división).
\end{formula}

\subsection{Índice de Gini}

Probabilidad de clasificar incorrectamente una muestra aleatoria etiquetada según la distribución
del nodo. Más eficiente computacionalmente que la entropía (sin logaritmos).

\begin{formula}[— Impureza de Gini]
  \[
    \text{Gini}(S) = 1 - \sum_i p_i^2
  \]
  \[
    \text{Gini\_split} = \frac{|L|}{|L|+|R|}\cdot\text{Gini}(L)
                       + \frac{|R|}{|L|+|R|}\cdot\text{Gini}(R)
  \]
\end{formula}

\subsubsection{Ejemplo: nodo con 4 elementos (3A, 1B)}
\[
  p_A = 0.75,\quad p_B = 0.25
\]
\[
  \text{Gini} = 1 - (0.75^2 + 0.25^2) = 1 - (0.5625 + 0.0625) = \mathbf{0.375}
\]
\textit{Interpretación:} 37.5\% de probabilidad de clasificar incorrectamente.

\begin{resumen}[— Entropía vs. Gini]
  \begin{tabular}{lll}
    \toprule
    \textbf{Criterio} & \textbf{Entropía} & \textbf{Gini}\\
    \midrule
    Cálculo & Requiere logaritmos (lento) & Solo potencias (rápido)\\
    Sensibilidad & Favorece nodos muy puros & Favorece divisiones equilibradas\\
    Resultado práctico & Casi idéntico & Idéntico en práctica\\
    Default sklearn & No & \textbf{Sí}\\
    \bottomrule
  \end{tabular}
\end{resumen}

\begin{cuando}[— Entropía vs. Gini]
  Usa \textbf{entropía} cuando quieras el árbol más informativo posible (diferencias muy pequeñas).
  Usa \textbf{Gini} cuando priorices la velocidad computacional (datasets grandes).
  En la práctica: no importa cuál elijas, el resultado será prácticamente idéntico.
\end{cuando}

\subsection{Criterios para Regresión}

\begin{center}
  \begin{tabular}{llp{5cm}}
    \toprule
    \textbf{Criterio} & \textbf{Fórmula} & \textbf{Cuándo usarlo}\\
    \midrule
    MSE (default) & $\frac{1}{n}\sum(y_i - \hat{y})^2$ & Errores grandes deben penalizarse más\\
    MAE & $\frac{1}{n}\sum|y_i - \hat{y}|$ & Hay outliers que no deben dominar\\
    Friedman MSE & MSE con corrección & Boosting (evalúa potencial de mejora)\\
    Poisson & $-\sum y_i \log\hat{y}_i + \hat{y}_i$ & Variables de conteo (eventos/tiempo)\\
    \bottomrule
  \end{tabular}
\end{center}

\section{Ventajas y Desventajas}

\begin{minipage}[t]{0.48\textwidth}
  \begin{tcolorbox}[colback=successbg, colframe=treegreen, title={\faCheck\ Ventajas}]
    \begin{itemize}[noitemsep]
      \item Fácil interpretación (caja blanca)
      \item Captura relaciones no lineales
      \item No requiere escalado de variables
      \item No requiere manejo de nulos
      \item Selección implícita de características
      \item Importancia de características disponible
    \end{itemize}
  \end{tcolorbox}
\end{minipage}
\hfill
\begin{minipage}[t]{0.48\textwidth}
  \begin{tcolorbox}[colback=critbg, colframe=accent, title={\faTimes\ Desventajas}]
    \begin{itemize}[noitemsep]
      \item Tendencia al sobreajuste sin regularización
      \item Alta varianza (sensible a datos)
      \item Predicciones discontinuas en regresión
      \item No los más precisos en solitario
      \item Algoritmo greedy $\Rightarrow$ no óptimo global
    \end{itemize}
  \end{tcolorbox}
\end{minipage}

\section{Regularización y Podado}

\begin{critico}
  Un árbol \textbf{sin regularización siempre sobreajusta}. Es la primera cosa que debes verificar.
  Train accuracy = 100\% + Test accuracy = 77\% = \textbf{señal de alarma}.
\end{critico}

\subsection{Pre-poda (durante la construcción)}

\begin{lstlisting}[caption={Pre-poda con GridSearchCV}]
from sklearn.tree import DecisionTreeClassifier
from sklearn.model_selection import GridSearchCV
import numpy as np

grid = {
    'max_depth':         [2, 3, 4, 5],       # Profundidad maxima
    'min_samples_leaf':  [3, 5, 7],           # Min muestras por hoja
    'min_samples_split': [2, 5, 10],          # Min muestras para dividir
}
gs = GridSearchCV(
    DecisionTreeClassifier(random_state=1),
    param_grid=grid,
    scoring='f1_weighted',
    cv=5
)
gs.fit(X_train, y_train)
print(f"Mejor: {gs.best_params_}")
\end{lstlisting}

\subsection{Post-poda: \texttt{ccp\_alpha}}

\begin{formula}[— Cost-Complexity Pruning]
  El parámetro \texttt{ccp\_alpha} ($\alpha$) penaliza la complejidad del árbol.
  A mayor $\alpha$, más ramas se podan. El algoritmo elimina primero los
  ``enlaces más débiles'' (ramas que menos aportan).
\end{formula}

\begin{lstlisting}[caption={Post-poda con ccp\_alpha}]
grid = {'ccp_alpha': np.logspace(-3, 3)}  # [0.001 ... 1000]
gs = GridSearchCV(
    DecisionTreeClassifier(random_state=1),
    param_grid=grid,
    scoring='recall'
)
gs.fit(X_train, y_train)
print(f"Mejor alpha: {gs.best_params_}")
# Resultado tipico: ccp_alpha=0.00543, recall=0.55
\end{lstlisting}

\begin{resumen}[— Pre-poda vs Post-poda]
  \begin{tabular}{lll}
    \toprule
    & \textbf{Pre-poda} & \textbf{Post-poda (ccp\_alpha)}\\
    \midrule
    Cuándo & Durante construcción & Tras construir árbol completo\\
    Velocidad & Más rápida & Más lenta\\
    Calidad & Puede ser subóptima & Generalmente más sistemática\\
    Uso recomendado & Datasets grandes & Cuando importa interpretabilidad\\
    \bottomrule
  \end{tabular}
\end{resumen}

\section{Invarianza al Escalado}

\begin{critico}
  Los árboles \textbf{NO se benefician del escalado de variables}.
  Sus decisiones se basan en umbrales; si escalas, el umbral se ajusta proporcionalmente
  y el resultado es \textbf{idéntico}. Aplicar \texttt{StandardScaler} a un árbol no daña,
  pero es tiempo perdido.
\end{critico}

\begin{cuando}[— Cuándo SÍ escalar con árboles]
  \begin{itemize}
    \item Cuando el pipeline incluye también modelos que requieren escalado (SVM, redes neuronales).
    \item Como paso previo a PCA o reducción de dimensionalidad antes del árbol.
    \item En algoritmos híbridos donde el árbol es un componente dentro de un pipeline complejo.
  \end{itemize}
\end{cuando}

\section{Guía de Ejercicios — Árboles}

\begin{ejercicio}[1: Cálculo manual de Entropía y Gini]
  Dado un nodo con 10 muestras: 6 de clase A y 4 de clase B.
  \begin{enumerate}
    \item Calcula la entropía $H(S)$.
    \item Calcula el índice de Gini.
    \item Tras dividir por un atributo: subconjunto izquierdo: 5A, 1B; derecho: 1A, 3B.
          Calcula la ganancia de información y el Gini split.
    \item ¿Cuál es el mejor criterio aquí? ¿Coinciden entropía y Gini en la decisión?
  \end{enumerate}
  \tcblower
  \textbf{Solución parcial:} $p_A=0.6$, $p_B=0.4$.
  $H(S) = -(0.6\log_2 0.6 + 0.4\log_2 0.4) = 0.971$.
  $\text{Gini}(S) = 1-(0.36+0.16)=0.48$.
\end{ejercicio}

\begin{ejercicio}[2: Diagnóstico de sobreajuste]
  Tienes un árbol con Train acc=0.98 y Test acc=0.74. Paso a paso:
  \begin{enumerate}
    \item ¿Qué tipo de error predomina: sesgo o varianza?
    \item ¿Qué estrategia de regularización aplicarías primero?
    \item Si tras \texttt{max\_depth=4} obtienes Train=0.84, Test=0.82, ¿está resuelto?
    \item ¿Tiene sentido aplicar \texttt{StandardScaler} para mejorar el resultado?
  \end{enumerate}
\end{ejercicio}

% ════════════════════════════════════════════════════════════
\chapter{Métodos de Ensamble — Teoría}
% ════════════════════════════════════════════════════════════

\section{¿Qué son los Métodos de Ensamble?}

Los métodos de ensamble combinan múltiples \textit{clasificadores débiles} (weak learners)
para formar un predictor más robusto. Un clasificador débil solo acierta ligeramente mejor
que el azar.

\begin{formula}[— Predictor de Ensamble]
  \[
    \hat{f}_{\text{ens}}(x) = \sum_{m=1}^{M} \alpha_m \cdot \hat{f}_m(x)
  \]
  donde $\hat{f}_m(x)$ es el $m$-ésimo estimador entrenado en $D_m \subseteq D$.
\end{formula}

\section{El Trade-off Sesgo-Varianza}

\begin{formula}[— Descomposición Sesgo-Varianza-Ruido]
  \[
    \mathbb{E}\bigl[(y_i - \hat{f}(x_i))^2\bigr]
    = \underbrace{\bigl[f(x_i) - \mathbb{E}(\hat{f}(x_i))\bigr]^2}_{\text{Sesgo}^2}
    + \underbrace{\mathbb{E}\bigl[(\hat{f}(x_i) - \mathbb{E}(\hat{f}(x_i)))^2\bigr]}_{\text{Varianza}}
    + \underbrace{\mathbb{E}[\varepsilon^2]}_{\text{Ruido irreducible}}
  \]
\end{formula}

\begin{resumen}[— Componentes del error]
  \begin{tabular}{llll}
    \toprule
    \textbf{Componente} & \textbf{Definición} & \textbf{Causa} & \textbf{Solución}\\
    \midrule
    Sesgo & Error sistemático & Modelo demasiado simple & Boosting, modelo complejo\\
    Varianza & Sensibilidad a los datos & Overfitting & Bagging, regularización\\
    Ruido & Variabilidad inherente & Datos ruidosos & No reducible con el modelo\\
    \bottomrule
  \end{tabular}
\end{resumen}

\begin{tip}[— Analogía del arquero]
  \begin{itemize}
    \item \textbf{Alto sesgo, baja varianza:} flechas juntas pero lejos del centro.
    \item \textbf{Bajo sesgo, alta varianza:} flechas dispersas alrededor del centro.
    \item \textbf{Bajo sesgo, baja varianza:} todas cerca del centro. \textit{El ideal.}
    \item \textbf{Alto sesgo, alta varianza:} el peor caso.
  \end{itemize}
\end{tip}

\section{Taxonomía de Métodos de Ensamble}

\begin{center}
  \begin{tabular}{lll}
    \toprule
    \textbf{Característica} & \textbf{BAGGING} & \textbf{BOOSTING}\\
    \midrule
    Objetivo & Reducir \textbf{varianza} & Reducir \textbf{sesgo}\\
    Entrenamiento & Paralelo (independiente) & Secuencial (corrige al anterior)\\
    Datos & Bootstrap (con reemplazo) & Re-pesados o residuos\\
    Combinación & Promedio o voto mayoritario & Suma ponderada\\
    Riesgo & No reduce sesgo alto & Puede sobreajustarse\\
    Paralelizable & Sí & No (secuencial por naturaleza)\\
    Ejemplos & Bagging, Random Forest & AdaBoost, GB, XGBoost, LGBM, CatBoost\\
    \bottomrule
  \end{tabular}
\end{center}

\begin{cuando}[— Bagging vs Boosting]
  \begin{itemize}
    \item Usa \textbf{Bagging/RF} cuando el árbol base ya es bastante bueno pero inestable (alta varianza).
    \item Usa \textbf{Boosting} cuando el árbol base es débil (underfitting) y necesitas reducir el sesgo.
    \item En la práctica, \textbf{Boosting casi siempre supera a Bagging} en precisión, pero es más
          propenso al sobreajuste y requiere más tuning.
  \end{itemize}
\end{cuando}

% ════════════════════════════════════════════════════════════
\chapter{Bagging y Random Forest}
% ════════════════════════════════════════════════════════════

\section{Bootstrap — Remuestreo con Reemplazo}

El bootstrap crea $M$ conjuntos de datos distintos a partir del dataset original. Selecciona
$n$ muestras \textbf{con reemplazo} de un dataset de $n$ observaciones. Resultado:
aproximadamente \textbf{63.2\%} de observaciones únicas por muestra (el 36.8\% son duplicados).

\begin{formula}[— ¿Por qué 63.2\%?]
  La probabilidad de que una muestra específica \textbf{no} sea seleccionada en un único
  sorteo es $(1 - 1/n)$. Tras $n$ sorteos:
  \[
    P(\text{no seleccionado}) = \left(1 - \frac{1}{n}\right)^n
    \xrightarrow{n\to\infty} e^{-1} \approx 0.368
  \]
  Por tanto, el 63.2\% sí aparece en cada muestra bootstrap.
\end{formula}

\section{Bagging}

\begin{formula}[— Predictor Bagging]
  \[
    \hat{f}_{\text{bag}}(x) = \frac{1}{M} \sum_{m=1}^{M} \hat{f}_m(x)
  \]
  \[
    \text{var}\bigl(\hat{f}_{\text{bag}}(x)\bigr) = \frac{\sigma^2}{M} \to 0
    \quad \text{cuando } M \to \infty
  \]
  \textit{La varianza se divide entre el número de estimadores.}
\end{formula}

\subsection{Out-of-Bag (OOB) — Validación Gratuita}

El $\sim$36.8\% de datos no seleccionados en cada bootstrap sirve como conjunto de validación
natural para ese árbol.

\begin{critico}
  El \textbf{error OOB} es un estimador casi tan bueno como la validación cruzada $k$-fold,
  pero computacionalmente mucho más eficiente: no requiere re-entrenar el modelo.
  Siempre activa \texttt{oob\_score=True} en BaggingRegressor y RandomForest.
\end{critico}

\begin{lstlisting}[caption={BaggingRegressor con OOB}]
from sklearn.ensemble import BaggingRegressor
from sklearn.tree import DecisionTreeRegressor

bag_model = BaggingRegressor(
    estimator=DecisionTreeRegressor(max_depth=3),
    n_estimators=50,
    oob_score=True,   # Validacion OOB gratuita
    n_jobs=-1,        # Paralelizar
    random_state=1
)
bag_model.fit(X_train, y_train)
print(f'OOB Score (R2): {bag_model.oob_score_:.3f}')
\end{lstlisting}

\section{Random Forest}

Random Forest resuelve el problema de correlación entre árboles en Bagging: si existe una
característica muy dominante, todos los árboles la usarán en el nodo raíz y estarán
correlacionados entre sí, limitando la reducción de varianza.

\textbf{Solución:} en cada nodo se selecciona aleatoriamente un subconjunto de $m$ características
del total de $p$, y solo se busca la mejor división dentro de ese subconjunto.

\begin{formula}[— Algoritmo Random Forest]
  \begin{align*}
    \text{Regresión:} &\quad \hat{f}_{\text{rf}}(x) = \frac{1}{B}\sum_{b=1}^{B} \hat{f}_b(x)\\[6pt]
    \text{Clasificación:} &\quad \hat{f}_{\text{rf}}(x) = \operatorname*{argmax}_{k \in \mathcal{Y}}
      \sum_{b=1}^{B} \mathbf{1}[\hat{f}_b(x) = k]\\[6pt]
    \text{Valores típicos:} &\quad m = \sqrt{p} \text{ (clasificación)}, \quad
      m = p/3 \text{ (regresión)}
  \end{align*}
\end{formula}

\begin{resumen}[— Hiperparámetros clave de Random Forest]
  \begin{tabular}{llll}
    \toprule
    \textbf{Parámetro} & \textbf{Descripción} & \textbf{Efecto} & \textbf{Valor típico}\\
    \midrule
    \texttt{n\_estimators} & Número de árboles & Más $\to$ menos varianza & 100--500\\
    \texttt{max\_features} & Features por nodo ($m$) & Menor $\to$ menos correlación & \texttt{sqrt} / \texttt{p/3}\\
    \texttt{max\_depth} & Profundidad por árbol & Mayor $\to$ menos sesgo & \texttt{None}\\
    \texttt{min\_samples\_leaf} & Muestras mínimas/hoja & Mayor $\to$ regularización & 1\\
    \texttt{ccp\_alpha} & Post-poda por árbol & Mayor $\to$ árboles más simples & 0\\
    \texttt{oob\_score} & Validación OOB & Ahorra conjunto de validación & \texttt{True}\\
    \texttt{n\_jobs} & Paralelismo & -1 = todos los núcleos & -1\\
    \bottomrule
  \end{tabular}
\end{resumen}

\begin{alerta}[— Error común con n\_estimators]
  Aumentar \texttt{n\_estimators} indefinidamente no mejora linealmente. Más de 200--500 árboles
  rara vez mejora el desempeño pero sí aumenta el tiempo de predicción.
  Usa la curva de aprendizaje para encontrar el punto de rendimientos decrecientes.
\end{alerta}

\begin{lstlisting}[caption={Random Forest con RandomizedSearchCV}]
from sklearn.ensemble import RandomForestRegressor
from sklearn.model_selection import RandomizedSearchCV
from scipy.stats import loguniform

dist = {
    'n_estimators': range(1, 201),
    'ccp_alpha':    loguniform(1e-4, 100),
    'max_features': [4, 5, 6, 7],
}
rs = RandomizedSearchCV(
    RandomForestRegressor(random_state=1),
    param_distributions=dist,
    scoring='neg_root_mean_squared_error',
    n_iter=30, random_state=1
)
rs.fit(X_train, y_train)
print(rs.best_params_)
\end{lstlisting}

\section{Guía de Ejercicios — Bagging y RF}

\begin{ejercicio}[3: Bootstrap manual]
  Dado el dataset $\{1.0, 2.4, 5.3, 3.7, 4.1\}$ (5 elementos):
  \begin{enumerate}
    \item Genera a mano 3 muestras bootstrap de tamaño 5 (usando una tabla de números aleatorios
          o semilla fija).
    \item ¿Cuántas muestras únicas hay en promedio? ¿Concuerda con el 63.2\%?
    \item ¿Cuáles serían los elementos OOB de cada muestra?
  \end{enumerate}
\end{ejercicio}

\begin{ejercicio}[4: Decorrelación en Random Forest]
  Explica con tus propias palabras:
  \begin{enumerate}
    \item ¿Por qué Bagging puro NO reduce perfectamente la varianza aunque $M \to \infty$?
    \item ¿Cómo resuelve Random Forest este problema?
    \item Si tienes 10 características ($p=10$), ¿qué valor de $m$ usarías en clasificación?
          ¿Y en regresión?
    \item Si reduces $m$ de 10 a 2, ¿qué pasa con el sesgo? ¿Y con la varianza?
  \end{enumerate}
\end{ejercicio}

% ════════════════════════════════════════════════════════════
\chapter{AdaBoost}
% ════════════════════════════════════════════════════════════

\section{Concepto e Historia}

AdaBoost (Adaptive Boosting) fue propuesto por Freund \& Schapire en 1996 y recibió el premio
Gödel (equivalente al Nobel en Ciencias de la Computación Teórica). Demostró que podían
combinarse clasificadores débiles para formar uno fuerte.

\begin{resumen}[— Ficha técnica de AdaBoost]
  \begin{tabular}{ll}
    \toprule
    \textbf{Origen} & Freund \& Schapire, 1996 (Premio Gödel)\\
    \textbf{Idea central} & Re-pesar muestras mal clasificadas\\
    \textbf{Clasificador base} & Stump (árbol de profundidad 1)\\
    \textbf{Predicción final} & $\hat{y} = \text{sign}\!\left(\sum \alpha_t h_t(x)\right)$\\
    \textbf{Función de pérdida} & Pérdida exponencial\\
    \textbf{Limitaciones} & Sensible a ruido; no soporta NaN\\
    \bottomrule
  \end{tabular}
\end{resumen}

\section{Algoritmo AdaBoost Paso a Paso}

\begin{formula}[— Algoritmo AdaBoost Clasificación Binaria]
  \textbf{Inicialización:} $w_i^{(1)} = \frac{1}{N}$ para $i = 1, \ldots, N$

  \medskip
  \textbf{Para $t = 1, \ldots, T$:}
  \begin{enumerate}
    \item Entrenar $h_t(x)$ con pesos $\mathbf{w}^{(t)}$.
    \item Calcular el error ponderado:
    \[
      \varepsilon_t = \frac{\sum_i w_i \cdot \mathbf{1}[y_i \neq h_t(x_i)]}{\sum_i w_i}
    \]
    \item Calcular el peso del estimador:
    \[
      \alpha_t = \frac{1}{2} \ln\!\left(\frac{1 - \varepsilon_t}{\varepsilon_t}\right)
    \]
    \item Actualizar pesos de muestras:
    \[
      w_i^{(t+1)} = w_i^{(t)} \cdot \exp\!\left(-\alpha_t \cdot y_i \cdot h_t(x_i)\right)
    \]
    Mal clasificada: $w$ \textbf{aumenta}. Bien clasificada: $w$ disminuye.
    \item Normalizar: $w_i^{(t+1)} \leftarrow w_i^{(t+1)} / \sum_j w_j^{(t+1)}$
  \end{enumerate}

  \medskip
  \textbf{Predicción final:} $\hat{y}(x) = \text{sign}\!\left(\sum_{t=1}^{T} \alpha_t h_t(x)\right)$
\end{formula}

\begin{tip}[— Intuición de AdaBoost]
  \begin{itemize}
    \item El primer clasificador intenta resolver el problema completo con pesos iguales.
    \item Las muestras fallidas se marcan como ``difíciles'' y reciben más atención.
    \item Cada siguiente clasificador se especializa en esas muestras difíciles.
    \item El consenso ponderado final es más robusto que cualquiera individualmente.
  \end{itemize}
\end{tip}

\begin{critico}
  \textbf{AdaBoost no tolera NaN ni variables categóricas directamente.}
  Siempre aplica \texttt{dropna()} y codifica categóricas antes de usarlo.
\end{critico}

\section{Implementación}

\begin{lstlisting}[caption={AdaBoost con Pipeline completo}]
from sklearn.ensemble import AdaBoostClassifier
from sklearn.pipeline import Pipeline
from sklearn.model_selection import RandomizedSearchCV
from scipy.stats import loguniform

ada = AdaBoostClassifier(n_estimators=50, random_state=1)
model = Pipeline([('preprocessor', preprocessor), ('classifier', ada)])
model.fit(X_train, y_train)

# Sintonizacion de hiperparametros
dist = {
    'classifier__learning_rate': loguniform(1e-3, 10),
    'classifier__n_estimators':  range(1, 200)
}
rs = RandomizedSearchCV(
    model, dist, scoring='f1_weighted',
    n_iter=20, random_state=1
)
rs.fit(X_train, y_train)
# Best: learning_rate=1.63, n_estimators=62
\end{lstlisting}

\section{Guía de Ejercicios — AdaBoost}

\begin{ejercicio}[5: Seguimiento manual de AdaBoost]
  Tienes 6 muestras con etiquetas $y = [+1, +1, +1, -1, -1, -1]$.
  El clasificador $h_1$ predice $[+1, +1, -1, -1, -1, +1]$.
  \begin{enumerate}
    \item ¿Cuáles muestras clasifica mal $h_1$?
    \item Calcula $\varepsilon_1$ con pesos iniciales iguales.
    \item Calcula $\alpha_1$.
    \item Actualiza los pesos de las muestras mal clasificadas vs bien clasificadas.
    \item Normaliza los pesos para que sumen 1.
  \end{enumerate}
\end{ejercicio}

% ════════════════════════════════════════════════════════════
\chapter{Gradient Boosting}
% ════════════════════════════════════════════════════════════

\section{La Idea Central}

Gradient Boosting reformula el boosting como un \textbf{problema de optimización}: en lugar de
re-pesar muestras, entrena cada nuevo árbol para predecir el \textbf{gradiente negativo} de la
función de pérdida respecto a las predicciones actuales — es decir, los \textbf{residuos}.

\begin{formula}[— Gradient Boosting: actualización]
  \[
    F_m(x) = F_{m-1}(x) + \eta \cdot h_m(x)
  \]
  \[
    F_M(x) = F_0(x) + \eta \cdot h_1(x) + \eta \cdot h_2(x) + \cdots + \eta \cdot h_M(x)
  \]
  donde $\eta$ es el \textit{learning rate} y $h_m$ es el $m$-ésimo árbol que predice los residuos.
\end{formula}

\section{Algoritmo — Clasificación Binaria}

\begin{formula}[— Gradient Boosting: clasificación]
  \begin{enumerate}
    \item \textbf{Inicialización:} $F_0(x) = \log\!\left(\frac{P(y=1)}{P(y=0)}\right)$
    \item \textbf{Para $m = 1, \ldots, M$:}
    \begin{enumerate}
      \item Residuos (gradiente negativo de la log-loss):
            $r_{im} = y_i - \hat{p}_{m-1}(x_i)$ donde
            $\hat{p}_{m-1} = \sigma(F_{m-1}(x_i)) = \frac{1}{1+e^{-F_{m-1}(x_i)}}$
      \item Entrenar $h_m(x)$ (árbol de regresión) para predecir $\{r_{im}\}$
      \item Actualizar: $F_m(x) = F_{m-1}(x) + \eta \cdot h_m(x)$
    \end{enumerate}
    \item \textbf{Predicción final:} $\hat{p}(x) = \sigma(F_M(x))$
  \end{enumerate}
  \medskip
  \textbf{Para regresión (MSE):} $F_0(x) = \bar{y}$; $r_{im} = y_i - F_{m-1}(x_i)$;
  $\hat{y}(x) = F_M(x)$.
\end{formula}

\begin{resumen}[— Hiperparámetros clave de Gradient Boosting]
  \begin{tabular}{lllll}
    \toprule
    \textbf{Parámetro} & \textbf{Descripción} & \textbf{Impacto} & \textbf{Rango típico}\\
    \midrule
    \texttt{n\_estimators} & Número de árboles & + $\to$ - sesgo, riesgo overfitting & 50--500\\
    \texttt{learning\_rate} & Escala por árbol & - $\to$ necesita más árboles & 0.01--0.3\\
    \texttt{max\_depth} & Profundidad por árbol & + $\to$ + varianza & 3--6\\
    \texttt{subsample} & Fracción datos/árbol & - $\to$ más regularización & 0.5--1.0\\
    \texttt{n\_iter\_no\_change} & Early stopping & Previene overfitting & 5--15\\
    \bottomrule
  \end{tabular}
\end{resumen}

\begin{tip}[— Regla práctica learning rate vs n\_estimators]
  Empieza con \texttt{lr=0.1} y \texttt{n\_estimators=100}. Si quieres afinar:
  reduce \texttt{lr} a 0.05 o 0.01 y multiplica \texttt{n\_estimators} proporcionalmente.
  Valida con early stopping.
\end{tip}

\section{Guía de Ejercicios — Gradient Boosting}

\begin{ejercicio}[6: Trazado de curva de aprendizaje]
  Con \texttt{GradientBoostingClassifier} y el dataset Adult:
  \begin{enumerate}
    \item Entrena con \texttt{learning\_rate=0.1} y \texttt{n\_estimators} en $\{10, 50, 100, 200, 500\}$.
    \item Registra el F1 de train y test para cada valor.
    \item ¿En qué punto empieza el overfitting?
    \item Repite con \texttt{learning\_rate=0.01}. ¿Cambia el punto de overfitting?
    \item ¿Qué concluyes sobre la relación entre \texttt{lr} y \texttt{n\_estimators}?
  \end{enumerate}
\end{ejercicio}

% ════════════════════════════════════════════════════════════
\chapter{Boosting Avanzado: XGBoost, LightGBM y CatBoost}
% ════════════════════════════════════════════════════════════

\section{XGBoost}

\subsection{Mejoras sobre Gradient Boosting clásico}

\begin{resumen}[— Innovaciones de XGBoost]
  \begin{tabular}{lp{5.5cm}p{4.5cm}}
    \toprule
    \textbf{Mejora} & \textbf{Problema que resuelve} & \textbf{Cómo funciona}\\
    \midrule
    Regularización ($\gamma$, $\lambda$) & GB puede sobreajustarse & $\gamma$ penaliza hojas extra;
      $\lambda$ es L2 sobre pesos\\
    Newton Boosting & GB usa solo 1er orden (lento) & Usa el Hessiano (2ª derivada);
      converge más rápido\\
    Cuantiles ponderados & Evaluar todos los splits es $O(n{\times}p)$ & Agrupa en bins ponderados
      por el Hessiano; $O(\text{bins}{\times}p)$\\
    NaN nativo & Requería imputación previa & Aprende la dirección óptima para valores faltantes\\
    Categorías nativas & One-Hot crea matrices dispersas & Ordena por estadísticas de gradiente\\
    \bottomrule
  \end{tabular}
\end{resumen}

\begin{lstlisting}[caption={XGBoost con early stopping y categorías nativas}]
from xgboost import XGBClassifier
from sklearn.model_selection import train_test_split

# Convertir categoricas a tipo 'category'
X[cat_cols] = X[cat_cols].astype('category')

X_tr, X_val, y_tr, y_val = train_test_split(
    X_train, y_train, test_size=0.1, stratify=y_train
)

model = XGBClassifier(
    grow_policy='lossguide',    # Leaf-wise (como LightGBM)
    tree_method='hist',         # Histogramas (rapido)
    enable_categorical=True,    # Soporte nativo categoricas
    early_stopping_rounds=10,
    random_state=1
)
model.fit(X_tr, y_tr, eval_set=[(X_val, y_val)], verbose=False)
\end{lstlisting}

\section{LightGBM}

\subsection{Las 4 Innovaciones}

\begin{resumen}[— Innovaciones de LightGBM]
  \begin{tabular}{lp{5.5cm}p{4.5cm}}
    \toprule
    \textbf{Innovación} & \textbf{Descripción} & \textbf{Beneficio}\\
    \midrule
    Histogramas & Discretiza features en bins ($\approx$255); busca splits en bordes &
      40,000$\times$ más rápido con 10M filas\\
    Leaf-wise growth & Solo expande la hoja con mayor ganancia & Mismo número de nodos,
      menor error\\
    GOSS & Usa muestras con gradiente alto; muestra aleatoria del resto &
      Reduce datos por árbol sin perder precisión\\
    EFB & Fusiona features que nunca son nonzero simultáneamente &
      Reduce dimensionalidad en datos dispersos\\
    \bottomrule
  \end{tabular}
\end{resumen}

\begin{alerta}[— Leaf-wise y overfitting]
  Leaf-wise puede crear árboles muy desequilibrados y sobreajustar.
  Controla la complejidad con \texttt{num\_leaves} (no solo con \texttt{max\_depth}).
  Regla: \texttt{num\_leaves} $< 2^{\texttt{max\_depth}}$.
\end{alerta}

\begin{lstlisting}[caption={LightGBM con early stopping}]
import lightgbm as lgb
from lightgbm import LGBMClassifier

lgb_model = LGBMClassifier(objective='binary', verbose=0)
lgb_model.fit(
    X_tr, y_tr,
    eval_set=[(X_val, y_val)],
    eval_metric='logloss',
    callbacks=[lgb.early_stopping(stopping_rounds=10)],
)
\end{lstlisting}

\section{CatBoost}

\subsection{Innovaciones}

\begin{resumen}[— Innovaciones de CatBoost]
  \begin{tabular}{lp{5cm}p{5cm}}
    \toprule
    \textbf{Innovación} & \textbf{Problema} & \textbf{Solución}\\
    \midrule
    Ordered Target Statistics & Target encoding clásico introduce fuga de información &
      Calcula la media del target usando \textbf{solo} las muestras anteriores al orden aleatorio\\
    Feature combinations & Interacciones entre categóricas & Crea combinaciones automáticas\\
    Árboles simétricos & Predicción costosa con árboles asimétricos &
      Misma condición en todos los nodos de un nivel; predicción $O(1)$\\
    \bottomrule
  \end{tabular}
\end{resumen}

\begin{critico}
  \textbf{CatBoost soporta NaN en numéricas, pero NO en categóricas.}
  Aplica \texttt{dropna(subset=categorical\_features)} antes de entrenar.
\end{critico}

\begin{lstlisting}[caption={CatBoost con soporte nativo de categóricas}]
from catboost import CatBoostClassifier

cat_model = CatBoostClassifier(verbose=0)
cat_model.fit(
    X_tr, y_tr,
    cat_features=categorical_features,   # Indicar categoricas
    eval_set=(X_val, y_val),
    early_stopping_rounds=10
)
\end{lstlisting}

\section{Comparativa de los 3 Algoritmos Avanzados}

\begin{center}
  \begin{tabular}{llll}
    \toprule
    \textbf{Criterio} & \textbf{XGBoost} & \textbf{LightGBM} & \textbf{CatBoost}\\
    \midrule
    Velocidad & Media & Muy alta & Media-alta\\
    Datos $>$1M filas & Lento & Excelente & Medio\\
    Categorías & Encoding manual & Básico nativo & Avanzado nativo (sin leakage)\\
    Valores nulos & Sí (nativo) & Sí (nativo) & Solo numéricas\\
    Regularización & $\gamma$, $\lambda$, $\alpha$ & \texttt{reg\_alpha/lambda} & \texttt{l2\_leaf\_reg}\\
    Crecimiento & Level-wise (configurable) & Leaf-wise (default) & Simétrico\\
    Mejor para & Datos mixtos, tuning fino & Datasets muy grandes & Muchas categóricas\\
    \bottomrule
  \end{tabular}
\end{center}

% ════════════════════════════════════════════════════════════
\chapter{Conceptos Fundamentales — Los 15 Más Críticos}
% ════════════════════════════════════════════════════════════

\begin{enumerate}[leftmargin=*, label=\textbf{\arabic*.}]

\item \textbf{Sobreajuste (Overfitting)}\\
  Train acc=100\%, test acc=77\% = señal de alarma. El árbol memorizó el ruido.
  \textit{Siempre regularizar.}

\item \textbf{Entropía y Ganancia de Información}\\
  La entropía mide cuánta impureza hay; la ganancia mide cuánto la reduce una división.
  \textit{Corazón del árbol: sin esto, no sabe qué pregunta hacer primero.}

\item \textbf{Impureza de Gini}\\
  Probabilidad de clasificar mal una muestra aleatoria. Más eficiente que entropía.
  \textit{En práctica, entropía y Gini dan resultados casi idénticos.}

\item \textbf{Podado (Pruning)}\\
  Pre-poda: más eficiente. Post-poda con \texttt{ccp\_alpha}: más sistemática.
  \textit{Siempre preferir post-poda cuando el tiempo lo permite.}

\item \textbf{Feature Importance}\\
  Importancia relativa basada en reducción de impureza. Herramienta de selección.
  \textit{No indica causalidad. Solo relevancia predictiva.}

\item \textbf{Bootstrap}\\
  $\approx$63.2\% de muestras únicas por muestra bootstrap. Permite crear $M$ datasets distintos.
  \textit{Fundamento de todo Bagging y Random Forest.}

\item \textbf{Out-of-Bag (OOB) Error}\\
  El 36.8\% no seleccionado sirve como validación gratuita por árbol.
  \textit{OOB $\approx$ CV $k$-fold, pero mucho más eficiente computacionalmente.}

\item \textbf{Decorrelación (Random Forest)}\\
  RF selecciona $m < p$ features por nodo, reduciendo correlación entre árboles.
  \textit{Sin decorrelación, el promedio de árboles similares mejora poco la varianza.}

\item \textbf{Sesgo-Varianza Trade-off}\\
  Error = Sesgo$^2$ + Varianza + Ruido. Bagging $\to$ reduce varianza. Boosting $\to$ reduce sesgo.
  \textit{Este tradeoff determina qué familia de algoritmo usar.}

\item \textbf{Learning Rate ($\eta$)}\\
  Escala la contribución de cada árbol. Menor $\eta$ $\Rightarrow$ más árboles, mejor generalización.
  \textit{Regla: lr=0.1, 100 árboles como punto de partida.}

\item \textbf{Early Stopping}\\
  Detener cuando el error de validación deja de mejorar. Previene overfitting automáticamente.
  \textit{Disponible en XGBoost, LightGBM y CatBoost. Siempre activarlo.}

\item \textbf{Target Encoding}\\
  Reemplaza cada categoría con la media del target. Útil para alta cardinalidad.
  \textit{Siempre usar TargetEncoder de sklearn (cross-fitting interno) para evitar leakage.}

\item \textbf{Residuos en Gradient Boosting}\\
  Los residuos = gradiente negativo de la pérdida. Cada árbol corrige los errores actuales.
  \textit{Clave: GB predice el ERROR, no $y$ directamente.}

\item \textbf{Hessiano (Newton Boosting — XGBoost)}\\
  La segunda derivada de la pérdida permite una aproximación cuadrática más precisa.
  \textit{XGBoost converge con menos árboles que GB clásico.}

\item \textbf{Stacking vs Averaging}\\
  Averaging: promedio simple. Stacking: meta-modelo que aprende cómo combinar modelos base.
  \textit{En competencias, stacking multi-nivel da típicamente el mejor resultado.}

\end{enumerate}

% ════════════════════════════════════════════════════════════
\chapter{Cheatsheet y Fórmulas de Referencia}
% ════════════════════════════════════════════════════════════

\section{Tabla de Fórmulas Esenciales}

\begin{longtable}{lll}
  \toprule
  \textbf{Concepto} & \textbf{Fórmula} & \textbf{Rango / Interpretación}\\
  \midrule
  \endhead
  Entropía & $H(S) = -\sum_c p(c)\log_2 p(c)$ & $[0, \log_2 k]$; 0=puro\\
  Ganancia info. & $GI(S,a) = H(S) - \sum_v \frac{|S_v|}{|S|}H(S_v)$ & $[0, H(S)]$; mayor=mejor split\\
  Gini & $\text{Gini}(S) = 1-\sum_i p_i^2$ & $[0, 1-1/k]$; 0=puro\\
  Gini split & $\frac{|L|}{|L|+|R|}\text{Gini}(L)+\frac{|R|}{|L|+|R|}\text{Gini}(R)$ & Minimizar\\
  Bagging & $\hat{f}_\text{bag}(x) = \frac{1}{M}\sum \hat{f}_m(x)$ & Promedio $M$ modelos\\
  Var. bagging & $\text{var}(\hat{f}_\text{bag}) = \sigma^2/M \to 0$ & Decrece con $M$\\
  RF regresión & $\hat{f}_\text{rf}(x) = \frac{1}{B}\sum_b \hat{f}_b(x)$ & Promedio de $B$ árboles\\
  RF clasificación & $\hat{f}_\text{rf}(x) = \operatorname{argmax}_k\sum_b \mathbf{1}[\hat{f}_b(x)=k]$ & Voto mayoritario\\
  AdaBoost error & $\varepsilon_t = \sum_i w_i \mathbf{1}[y_i\neq h_t(x_i)] / \sum_i w_i$ & $[0, 0.5]$; $<0.5$ = útil\\
  AdaBoost $\alpha$ & $\alpha_t = \frac{1}{2}\ln\!\left(\frac{1-\varepsilon_t}{\varepsilon_t}\right)$ & Mayor $\alpha_t$ = más influencia\\
  AdaBoost pesos & $w_i^{(t+1)} = w_i^{(t)}\exp(-\alpha_t y_i h_t(x_i))$ & Aumenta pesos de errores\\
  AdaBoost pred. & $\hat{y}(x) = \text{sign}(\sum_t \alpha_t h_t(x))$ & Votación ponderada\\
  GB residuos & $r_{im} = -\partial L(y_i,F(x_i))/\partial F(x_i)$ & Gradiente negativo\\
  GB (MSE) & $r_{im} = y_i - F_{m-1}(x_i)$ & Diferencia directa\\
  GB actualiz. & $F_m(x) = F_{m-1}(x) + \eta h_m(x)$ & $\eta$=lr, $h_m$=árbol nuevo\\
  Sesgo-Varianza & $E_\text{total} = \text{Sesgo}^2 + \text{Varianza} + \text{Ruido}$ & Tradeoff fundamental\\
  Bootstrap & $P(\text{no sel.}) = (1-1/n)^n \approx e^{-1} \approx 0.368$ & 63.2\% únicos/muestra\\
  \bottomrule
  \caption{Fórmulas esenciales del módulo}
\end{longtable}

\section{Guía Rápida de Implementación}

\begin{lstlisting}[caption={Cheatsheet de código — todos los modelos}]
# === ARBOL DE DECISION ========================================
from sklearn.tree import DecisionTreeClassifier
clf = DecisionTreeClassifier(
    criterion='gini',         # 'gini' (default) o 'entropy'
    max_depth=5,              # Pre-poda
    min_samples_leaf=5,       # Pre-poda
    ccp_alpha=0.01,           # Post-poda
    random_state=1
)

# === RANDOM FOREST ============================================
from sklearn.ensemble import RandomForestClassifier
rf = RandomForestClassifier(
    n_estimators=200,
    max_features='sqrt',      # m = sqrt(p)
    oob_score=True,
    n_jobs=-1, random_state=1
)

# === ADABOOST =================================================
from sklearn.ensemble import AdaBoostClassifier
ada = AdaBoostClassifier(
    n_estimators=100,
    learning_rate=1.0,
    random_state=1
)

# === GRADIENT BOOSTING ========================================
from sklearn.ensemble import GradientBoostingClassifier
gb = GradientBoostingClassifier(
    n_estimators=100, learning_rate=0.1, max_depth=3,
    subsample=0.8, n_iter_no_change=10,
    validation_fraction=0.1, random_state=1
)

# === XGBOOST ==================================================
from xgboost import XGBClassifier
xgb = XGBClassifier(
    n_estimators=500, learning_rate=0.05, max_depth=6,
    subsample=0.8, colsample_bytree=0.8,
    gamma=1, reg_alpha=0.1, reg_lambda=1.0,
    scale_pos_weight=3,       # Desbalance: neg/pos
    tree_method='hist', enable_categorical=True,
    early_stopping_rounds=10, random_state=1
)

# === LIGHTGBM =================================================
import lightgbm as lgb
lgbm = lgb.LGBMClassifier(
    n_estimators=1000, learning_rate=0.05,
    num_leaves=31,            # Control leaf-wise
    min_child_samples=20,     # Regularizacion
    subsample=0.8, colsample_bytree=0.8,
    reg_alpha=0.1, reg_lambda=1.0, verbose=-1
)
lgbm.fit(X_tr, y_tr, eval_set=[(X_val, y_val)],
         callbacks=[lgb.early_stopping(50)])

# === CATBOOST =================================================
from catboost import CatBoostClassifier
cat = CatBoostClassifier(
    iterations=1000, learning_rate=0.05,
    depth=6, l2_leaf_reg=3, verbose=0
)
cat.fit(X_tr, y_tr, cat_features=['col1','col2'],
        eval_set=(X_val, y_val), early_stopping_rounds=50)

# === BUSQUEDA DE HIPERPARAMETROS ==============================
from sklearn.model_selection import RandomizedSearchCV
from scipy.stats import loguniform, uniform, randint
dist = {'learning_rate': loguniform(0.01, 0.3),
        'n_estimators':  randint(50, 500)}
rs = RandomizedSearchCV(model, dist, n_iter=50,
    scoring='f1_weighted', cv=5, n_jobs=-1, random_state=1)
rs.fit(X_train, y_train)
\end{lstlisting}

% ════════════════════════════════════════════════════════════
\chapter{Guía de Estudio Definitiva}
% ════════════════════════════════════════════════════════════

\section{Prerrequisitos}

\begin{center}
  \begin{tabular}{lll}
    \toprule
    \textbf{Prerrequisito} & \textbf{Nivel requerido} & \textbf{Por qué es necesario}\\
    \midrule
    Python (pandas, numpy, matplotlib) & Intermedio & Todos los notebooks lo usan\\
    Estadística básica & Básico & Probabilidad, media, varianza, distribuciones\\
    Álgebra lineal & Básico & Vectores y matrices para feature importances\\
    Teoría de la información & Nociones & Entropía y ganancia de información\\
    scikit-learn API & Básico-Intermedio & fit/predict/Pipeline/GridSearchCV\\
    Conceptos de ML & Básico & Overfitting, cross-validation, métricas\\
    \bottomrule
  \end{tabular}
\end{center}

\section{Roadmap de Aprendizaje — 4 Semanas}

\begin{tcolorbox}[colback=infobg, colframe=baggingblue, title={\faCalendar\ Semana 1 — Árboles de Decisión}]
  \begin{itemize}[noitemsep]
    \item Estudiar: capítulos 2 y 3 de esta guía (teoría CART, criterios de impureza)
    \item Implementar: árbol con Adult dataset; árbol de regresión con Auto-MPG
    \item Ejercicio: calcular manualmente entropía y Gini (ejercicios 1 y 2)
    \item Meta: entender CART, criterios de impureza, regularización
  \end{itemize}
\end{tcolorbox}

\begin{tcolorbox}[colback=infobg, colframe=baggingblue, title={\faCalendar\ Semana 2 — Bagging y Random Forest}]
  \begin{itemize}[noitemsep]
    \item Estudiar: capítulo 4 de esta guía (bootstrap, OOB, decorrelación)
    \item Implementar: BaggingRegressor vs RandomForestRegressor
    \item Ejercicio: bootstrap manual y análisis de decorrelación (ejercicios 3 y 4)
    \item Meta: entender bootstrap, OOB, $m$-features, Random Forest
  \end{itemize}
\end{tcolorbox}

\begin{tcolorbox}[colback=warningbg, colframe=boostingorange, title={\faCalendar\ Semana 3 — Boosting: AdaBoost y Gradient Boosting}]
  \begin{itemize}[noitemsep]
    \item Estudiar: capítulos 5 y 6 de esta guía
    \item Implementar: AdaBoost y GradientBoosting con Pipeline completo
    \item Ejercicio: seguimiento manual de AdaBoost (ejercicio 5); curva de aprendizaje (ejercicio 6)
    \item Meta: entender residuos, learning rate, funciones de pérdida
  \end{itemize}
\end{tcolorbox}

\begin{tcolorbox}[colback=critbg, colframe=accent, title={\faCalendar\ Semana 4 — XGBoost, LightGBM y CatBoost}]
  \begin{itemize}[noitemsep]
    \item Estudiar: capítulo 7 de esta guía + documentación oficial de cada librería
    \item Implementar: comparativa completa de los 3 en el mismo dataset
    \item Meta: saber cuándo usar cada uno, cómo configurar early stopping y tuning
  \end{itemize}
\end{tcolorbox}

\begin{tcolorbox}[colback=successbg, colframe=treegreen, title={\faTrophy\ Consolidación — Proyecto Integrador}]
  \begin{itemize}[noitemsep]
    \item Aplicar pipeline completo: EDA $\to$ preprocesamiento $\to$ árbol $\to$ RF $\to$ boosting
    \item Comparar modelos con tabla de resultados (F1 train/test, tiempo, overfitting)
    \item Analizar importancias de features con todos los modelos
    \item Implementar SHAP values para el mejor modelo
  \end{itemize}
\end{tcolorbox}

\section{Errores Comunes y Cómo Evitarlos}

\begin{center}
  \begin{tabular}{lp{4.5cm}p{5.5cm}}
    \toprule
    \textbf{Concepto} & \textbf{Error común} & \textbf{Estrategia correcta}\\
    \midrule
    \texttt{ccp\_alpha} & Confundirlo con umbral de precisión &
      Visualizar árbol con distintos $\alpha$; observar simplificación\\
    Residuos en GB & Creer que es $y - \hat{y}$ siempre &
      Para MSE sí; para log-loss es el gradiente negativo de la pérdida\\
    lr vs n\_estimators & Tratarlos como independientes &
      Son complementarios; lr bajo requiere más árboles. Usar early stopping\\
    OOB error & Confundirlo con error en test &
      OOB estima la generalización, no es exactamente igual a test\\
    Feature importance & Inferir causalidad &
      Solo indica relevancia predictiva. No causa-efecto\\
    Leaf-wise (LGBM) & Ignorar riesgo de overfitting &
      Controlar con \texttt{num\_leaves}, no solo \texttt{max\_depth}\\
    Target encoding & Aplicar sin protección &
      Siempre usar \texttt{TargetEncoder} de sklearn o CatBoost\\
    \bottomrule
  \end{tabular}
\end{center}

\section{Tabla Resumen Final}

\begin{center}
  \begin{tabular}{lll}
    \toprule
    \textbf{Necesidad} & \textbf{Algoritmo} & \textbf{Parámetro clave}\\
    \midrule
    Interpretabilidad total + reglas & Árbol podado (depth=3--5) & \texttt{ccp\_alpha} / \texttt{max\_depth}\\
    Baseline rápido y robusto & Random Forest & \texttt{max\_features='sqrt'}\\
    Precisión, datos medianos & Gradient Boosting / AdaBoost & \texttt{lr} + \texttt{n\_estimators}\\
    Mejor precisión general & XGBoost & \texttt{gamma}, \texttt{reg\_alpha/lambda}\\
    Datos muy grandes ($>$500k) & LightGBM & \texttt{num\_leaves}, \texttt{min\_child\_samples}\\
    Muchas variables categóricas & CatBoost & \texttt{cat\_features}, \texttt{l2\_leaf\_reg}\\
    Máxima precisión (competencia) & Stacking RF+XGB+LGBM+CAT & Meta-modelo logístico\\
    \bottomrule
  \end{tabular}
\end{center}

% ════════════════════════════════════════════════════════════
\chapter{Conclusiones y Próximos Pasos}
% ════════════════════════════════════════════════════════════

\section{Aprendizajes Principales}

\begin{enumerate}[label=\checkmark, leftmargin=*]
  \item \textbf{Los árboles siempre sobreajustan sin regularización.}
        Sin podado: train=100\%, test=77\%. La regularización es obligatoria.

  \item \textbf{Bagging reduce varianza; Boosting reduce sesgo.}
        Son filosofías complementarias. Si el árbol base es estable pero sesgado,
        Boosting es más útil. Si es inestable, Bagging primero.

  \item \textbf{Random Forest: robusto con poco tuning.}
        Solo ajustando \texttt{n\_estimators} y \texttt{max\_depth} ya da buenos resultados.
        Es el algoritmo ``seguro'' cuando no se sabe qué usar.

  \item \textbf{XGBoost/LightGBM/CatBoost son el estado del arte para tabular.}
        En prácticamente cualquier competencia con datos tabulares, uno de estos tres gana.
        En producción, LightGBM es favorito por velocidad.

  \item \textbf{El escalado NO ayuda a los árboles.}
        Invariantes al escalado por definición. No desperdiciar tiempo en \texttt{StandardScaler}.

  \item \textbf{Las categóricas requieren codificación (excepto CatBoost).}
        AdaBoost y GB clásico no soportan categóricas. XGBoost y LightGBM sí con configuración.

  \item \textbf{Los valores nulos no son problema (excepto AdaBoost).}
        XGBoost, LightGBM y CatBoost (en numéricas) manejan NaN nativamente.
\end{enumerate}

\section{Próximos Pasos Recomendados}

\begin{itemize}
  \item Implementar \textbf{SHAP values} para interpretabilidad post-hoc de ensambles
        (más informativo que \texttt{feature\_importances\_}).
  \item Explorar \textbf{Optuna} para búsqueda bayesiana de hiperparámetros
        (más eficiente que \texttt{RandomizedSearchCV}).
  \item Implementar \textbf{stacking}: árbol + RF + XGBoost + LGBM como base;
        meta-modelo logístico.
  \item Explorar \textbf{calibración de probabilidades} (Platt scaling, isotonic regression)
        para modelos de boosting.
  \item Aplicar \textbf{curvas de aprendizaje} para diagnosticar sesgo vs varianza antes de
        elegir el algoritmo.
  \item Investigar \textbf{pérdidas personalizadas} (focal loss para clases muy desbalanceadas).
\end{itemize}

\vfill
\begin{center}
  \textit{Guía Maestra generada el 21 de mayo de 2026}\\
  \textit{Emanuel Cabrera Novoa · Análisis Predictivo de Datos · Docente: Jhon Quiza}
\end{center}

\end{document}
